\section*{Chapitre \ref{ch:g10:habitable-planet} --- Une planète habitable}

\begin{solution}{exo:g10:habitable-planet:1}
La température à laquelle un émetteur parfait rayonnerait vers l'espace
exactement la puissance que la planète absorbe de son étoile. La surface de
la Terre est plus chaude (\qty{288}{K} contre environ \qty{255}{K}) parce
que les gaz à effet de serre absorbent une partie de l'infrarouge qu'elle
émet et en renvoient une partie vers le bas.
\end{solution}

\begin{solution}{exo:g10:habitable-planet:2}
Vénus~: \qty{737}{K} mesurés contre une \omterm{def:g10:habitable-planet:equilibrium-temperature}{température d'équilibre} d'environ
\qty{227}{K}, soit un réchauffement d'environ \qty{510}{K}.
\end{solution}

\begin{solution}{exo:g10:habitable-planet:3}
$S = 1361/1.52^2 \approx \qty{589}{W/m^2}$ (la fiche donne
\qty{586}{W/m^2}, avec la distance moyenne exacte).
\end{solution}

\begin{solution}{exo:g10:habitable-planet:4}
\qty{1}{bar}, \qty{330}{K}~: liquide (en dessous du point d'ébullition,
environ \qty{373}{K}). \qty{0.001}{bar}, \qty{250}{K}~: glace (la pression
est au-dessus de la courbe glace--vapeur, qui se trouve vers
\qty{0.0008}{bar} à \qty{250}{K}). \qty{100}{bar}, \qty{700}{K}~: vapeur
--- la température dépasse la température critique, donc aucun liquide ne
peut exister.
\end{solution}

\begin{solution}{exo:g10:habitable-planet:5}
Une planète qui tourne autour d'une autre étoile que le Soleil. Transits~:
la planète passe devant son étoile et en diminue légèrement l'éclat à
chaque orbite. Oscillation~: l'attraction de la planète fait aller et venir
l'étoile le long de notre ligne de visée, ce qui décale périodiquement la
couleur de sa lumière.
\end{solution}

\begin{solution}{exo:g10:habitable-planet:6}
$T = (1361 \times 0.706/(4 \times 5.67\times10^{-8}))^{1/4} \approx
\qty{255}{K} \approx \qty{-18}{\degreeCelsius}$. Avec $A = 0$~:
$\approx \qty{278}{K} \approx \qty{5}{\degreeCelsius}$.
\end{solution}

\begin{solution}{exo:g10:habitable-planet:7}
$\sqrt 4 = 2$~: de $2 \times 0.99 = \qty{1.98}{au}$ à $2 \times 1.70 =
\qty{3.40}{au}$. À \qty{1}{au}, la planète recevrait quatre fois
l'éclairement de la Terre~: bien en deçà du bord intérieur, trop chaude.
\end{solution}

\begin{solution}{exo:g10:habitable-planet:8}
$S = 1361/0.723^2 \approx \qty{2604}{W/m^2}$;
$T_{\mathrm{eq}} = (2604 \times 0.23/(4\sigma))^{1/4} \approx \qty{227}{K}$.
La surface est plus chaude d'environ \qty{510}{K}~: son épaisse atmosphère
de dioxyde de carbone (\qty{92}{bar}) produit un énorme effet de serre.
\end{solution}

\begin{solution}{exo:g10:habitable-planet:9}
À \qty{6.36}{mbar}, l'eau n'est liquide qu'entre environ
\qty{0}{\degreeCelsius} et \qty{0.5}{\degreeCelsius}. À
\qty{10}{\degreeCelsius}, la pression est en dessous de la courbe
d'ébullition~: l'eau s'évapore en bouillant~; la nuit, elle gèle.
\end{solution}

\begin{solution}{exo:g10:habitable-planet:10}
$11.2/2.4 \approx 4.7$. Une molécule de gaz doit aller presque cinq fois plus
vite pour échapper à la Terre qu'à la Lune~; en des milliards d'années, la
faible gravité de la Lune a laissé ses gaz s'échapper, tandis que celle de
la Terre a retenu l'essentiel des siens.
\end{solution}

\begin{solution}{exo:g10:habitable-planet:11}
$T = (1361 \times 0.50/(4\sigma))^{1/4} \approx \qty{234}{K}$, environ
\qty{-39}{\degreeCelsius}. La glace réfléchirait tant de lumière solaire que
la planète resterait assez froide pour la conserver~: une boucle qui se
renforce elle-même.
\end{solution}

\begin{solution}{exo:g10:habitable-planet:12}
$T_{\mathrm{eq}} = (S_0(1-A)/(4\sigma d^2))^{1/4} = (S_0/(4\sigma))^{1/4}
(1-A)^{1/4} (d^2)^{-1/4}$~; $(1361/(4\sigma))^{1/4} \approx \qty{278.3}{K}$ et
$(d^2)^{1/4} = \sqrt d$. Quatre fois plus loin~: $T$ est divisée par
$\sqrt 4 = 2$.
\end{solution}

\begin{solution}{exo:g10:habitable-planet:13}
$\sqrt d = 278.3 \times 0.706^{1/4}/273 \approx 0.934$, donc $d \approx
\qty{0.873}{au}$, en deçà du bord intérieur (\qty{0.99}{au}). La \omterm{def:g10:habitable-planet:equilibrium-temperature}{température
d'équilibre} ignore l'effet de serre~: une vraie planète dotée d'une
atmosphère est plus chaude que $T_{\mathrm{eq}}$, si bien qu'elle surchauffe
et perd son eau dès une distance plus grande, vers \qty{0.99}{au}. Les bords
viennent de modèles de l'atmosphère entière, pas de $T_{\mathrm{eq}}$ seule.
\end{solution}

\begin{solution}{exo:g10:habitable-planet:14}
Entre le point de fusion, environ \qty{0}{\degreeCelsius}, et le point
d'ébullition, environ \qty{7}{\degreeCelsius}. À \qty{15}{\degreeCelsius},
le lac bouillirait.
\end{solution}

\begin{solution}{exo:g10:habitable-planet:15}
$\sqrt{0.5} \approx 0.707$~: de $0.707 \times 0.99 \approx \qty{0.70}{au}$
à $0.707 \times 1.70 \approx \qty{1.20}{au}$. La planète est exactement au
bord extérieur. Les bords sont des résultats de modèles et la mise à
l'échelle est approchée~: impossible de dire si elle est à l'intérieur ou à
l'extérieur.
\end{solution}

\begin{solution}{pb:g10:habitable-planet:1}
\textbf{1.} $0.00151/0.0485^2 \approx 0.642$.

\textbf{2.} $0.642 \times 1361 \approx \qty{874}{W/m^2}$.

\textbf{3.} Entre les deux~: environ trois fois moins que Vénus, une fois et
demie plus que Mars.

\textbf{4.} $365.25/11.18 \approx 32.7$.

\textbf{5.} $(874/(4 \times 5.67\times10^{-8}))^{1/4} \approx \qty{249}{K}$.

\textbf{6.} $(874 \times 0.23/(4\sigma))^{1/4} \approx \qty{173}{K}$.

\textbf{7.} Aucune~: \qty{249}{K} reste inférieur à \qty{273}{K}. Une
couverture brillante de nuages ou de glace abaisse la \omterm{def:g10:habitable-planet:equilibrium-temperature}{température
d'équilibre} de plus de \qty{75}{K}~: ce qui recouvre une planète compte
autant que sa distance.

\textbf{8.} $\sqrt{0.00151} \approx 0.0389$.

\textbf{9.} $0.99 \times 0.0389 \approx \qty{0.0385}{au}$.

\textbf{10.} $1.70 \times 0.0389 \approx \qty{0.066}{au}$.

\textbf{11.} Oui~: $0.0385 < 0.0485 < 0.066$.

\textbf{12.} $d_\odot = 0.0485/0.0389 \approx \qty{1.25}{au}$~: dans la
moitié intérieure de la \omterm{def:g10:habitable-planet:habitable-zone}{zone habitable} du Soleil, entre la Terre et Mars.

\textbf{13.} $1.055^{1/3} \approx 1.018$.

\textbf{14.} $11.2 \times \sqrt{1.055/1.018} \approx \qty{11.4}{km/s}$, à
peu près celle de la Terre~: sa gravité pourrait retenir une atmosphère
(qu'elle en ait une dépend de son histoire).

\textbf{15.} L'oscillation de l'étoile ne révèle que la composante de son
mouvement le long de notre ligne de visée~; comme l'inclinaison de l'orbite
est inconnue, la masse trouvée n'est la vraie masse que si l'orbite est vue
par la tranche, et une valeur minimale sinon.

\textbf{16.} La pression du point triple, \qty{611.657}{Pa}, environ
\qty{0.006}{bar}.

\textbf{17.} $(874 \times 0.706/(4\sigma))^{1/4} \approx \qty{228}{K}$.

\textbf{18.} $288 - 228 = \qty{60}{K}$~: Proxima~b aurait besoin d'un
réchauffement par effet de serre \textbf{d'environ \qty{60}{K}}, presque le
double des \qty{33}{K} de la Terre.
\end{solution}
